\chapter{Microscopic Scale-Transfer Closure for the Vacuum Relaxation Rate}
\label{app:vacuumdamping_cosmicredshift}

\noindent\textbf{Abstract.}
This appendix develops a microscopic closure for the weak vacuum relaxation
rate $\Gamma_{\mathrm D}$ acting on the coherent photon branch of Tension
Field Theory (TFT). The construction addresses the microscopic origin of a
slow, coherent migration of a propagating photon through its real-frequency
eigenbranch. Ordinary linear amplitude damping does not generate such a
migration: it attenuates a fixed-wave-number carrier while leaving its real
frequency unchanged to leading order. The required TFT mechanism must instead
act on the self-similar photon scale coordinate.

The construction uses the established FCC transverse modal normalization
\begin{equation}
    A_{\mathrm{ct}}^{(\mathrm{modal})}
    =
    2\pi\ell_{\mathrm S}^{2},
\end{equation}
where $\ell_{\mathrm S}$ is the fixed Hard-Floor substrate spacing.
The quantity $A_{\mathrm{ct}}^{(\mathrm{modal})}$ remains a modal
surface-to-cell normalization and is not reinterpreted as a literal
microscopic contact patch.

A minimal reference-mediated scale-transfer closure assigns this transverse
modal support one substrate-reference correlation depth $\ell_{\mathrm S}$.
The resulting normalized reference overlap with one Soft-Floor node scales as
$(\ell_{\mathrm S}/r_{\mathrm D})^{3/2}$. Under the adopted
two-reference-vertex return closure, the coherent photon leaves and returns to
the Soft-Floor photon branch. The effective coupling is therefore second order
in the reference-overlap amplitude, and the corresponding scale-relaxation
kernel is suppressed by $(\ell_{\mathrm S}/r_{\mathrm D})^6$.

The resulting conditional microscopic rate is
\begin{equation}
    \boxed{
    \Gamma_{\mathrm D}
    =
    \frac{9}{16}\,
    \mathcal F_{\mathrm R}\,
    \frac{c_{\mathrm E}}{r_{\mathrm D}}
    \left(
        \frac{\ell_{\mathrm S}}{r_{\mathrm D}}
    \right)^6
    }
    \label{eq:AQ_gamma_master}
\end{equation}
where $\mathcal F_{\mathrm R}$ is a dimensionless
reference-mediated carrier--noncoherent correlation factor. It contains the
normalized intermediate-state response, the carrier-projected active-vacuum
spectral weight, and the physical correlation response relative to the
reference microscopic interval.

For the unit-reference closure
$\mathcal F_{\mathrm R}=1$, the presently calibrated TFT microscopic scales
give
\begin{equation}
    \Gamma_{\mathrm D}^{(0)}
    \simeq
    2.36\times10^{-18}\ {\rm s^{-1}},
\end{equation}
without using a Hubble constant, cosmic age, cosmological distance, CMB
temperature, supernova fit, or BAO scale as an upstream input.

The microscopic suppression factor is not interpreted as the probability of a
rare destructive photon collision. It is instead the small coherent weight of
a scale-projected response operating over successive microscopic correlation
cycles. Under preservation of the coherent modal action
$I_{\gamma}=E_{\gamma}/\omega_{\gamma}=\hbar$, this transfer produces
\begin{equation}
    \frac{d\ln\omega_{\gamma}}{dt}
    =
    -\Gamma_{\mathrm D}
\end{equation}
while preserving photon number, propagation direction, and polarization at
leading order within the selected coherent-channel closure. Energy transferred
out of the coherent photon branch is assigned only to noncoherent degrees of
freedom of the active vacuum at this stage; no thermal-equilibrium or CMB
identification is made in this appendix.

The resulting frequency shift defines only the relaxation-induced component
$z_{\mathrm D}$ of an astronomical redshift. TFT is globally non-expanding at
the Structural Reservoir level but locally dynamical in its matter and
phase-boundary sectors; therefore local source--observer kinematics may
generate an additional kinematic contribution $z_{\mathrm{kin}}$, while
different gravitational endpoint clock conditions contribute
$z_{\mathrm{grav}}$. In the stationary factorized regime,
\begin{equation}
    1+z_{\mathrm{obs}}
    =
    (1+z_{\mathrm{kin}})
    (1+z_{\mathrm{grav}})
    (1+z_{\mathrm D}),
\end{equation}
with the full observational decomposition developed separately. Consequently,
the characteristic length
$c_{\mathrm E}/\Gamma_{\mathrm D}$ is a photon scale-relaxation length and
must not be identified with the full Structural Reservoir radius
$R_{\mathrm{Res}}$.

The numerical value of Equation~\eqref{eq:AQ_gamma_master} remains sensitive
to the independent microscopic determination of $r_{\mathrm D}$. The present
TFT determination of that scale inherits an unresolved microscopic
two-dimensional-to-three-dimensional projection, and the numerical result
therefore remains calibration dependent at this stage.


\section{Scope, Architectural Setting, and Logical Status}

TFT distinguishes between the structural and kinematic descriptions of the
vacuum.

At the structural level, the theory is discrete and quantum. The Hard Floor
provides a phase-locked substrate reference, while the active FCC Soft Floor
contains finite-energy propagating degrees of freedom.

At the kinematic level, the long-wavelength dynamics are represented by a
continuum limit. Matter, localized phase structures, black-hole phase
boundaries, and other macroscopic structures may therefore move, rotate,
recede, approach, merge, or reorganize inside the active medium even when the
Structural Reservoir itself is not undergoing FLRW-type metric expansion.

Accordingly,
\begin{equation}
    \boxed{
    \text{global staticity}
    \;\not\Rightarrow\;
    \text{local kinematic staticity}.
    }
    \label{eq:AQ_global_local_distinction}
\end{equation}

This distinction is important for the interpretation of astronomical
observables but is not required for the microscopic derivation of
$\Gamma_{\mathrm D}$ itself.

The purpose of the present appendix is therefore narrow. In the deterministic
weak-transfer reduction, the physical relaxation rate is defined by
\begin{equation}
    \boxed{
    \Gamma_{\mathrm D}
    \equiv
    -\left.
    \frac{d\ln E_{\gamma}}{dt}
    \right|_{\mathrm{rel}},
    \qquad
    [\Gamma_{\mathrm D}]
    =
    \mathrm{s^{-1}}.
    }
\end{equation}
The qualifier $\mathrm{rel}$ denotes energy transferred from the coherent
photon branch into finite-energy Active-Vacuum receiving degrees of freedom.
It excludes source--observer Doppler shifts, gravitational endpoint clock
differences, and conservative frequency conversion by an evolving background.

It does not derive:

\begin{itemize}
    \item the global size of the Structural Reservoir,
    \item the macroscopic source velocity field,
    \item black-hole recession or rotational kinematics,
    \item the complete observed-redshift decomposition,
    \item the supernova light-curve timing law,
    \item or the BAO/CMB observational mapping.
\end{itemize}

Instead, it asks one specific question:

\begin{equation}
    \boxed{
    \begin{gathered}
        \text{Can the established TFT microscopic hierarchy support a finite}\\
        \Gamma_{\mathrm D}\quad
        \text{without cosmological calibration?}
    \end{gathered}
    }
\end{equation}

Three logical levels must be kept separate.

First, the following relations are inherited from the established TFT
construction:
\begin{align}
    d_{\mathrm{nn}}
    &=
    2r_{\mathrm D},
    \label{eq:AQ_dnn}
    \\
    A_{\mathrm{ct}}^{(\mathrm{modal})}
    &=
    \frac{M_{\mathrm D}}{\mu_{\mathrm D}}
    =
    2\pi\ell_{\mathrm S}^{2},
    \label{eq:AQ_Act}
    \\
    M_{\mathrm D}c_{\mathrm E}r_{\mathrm D}
    &=
    h.
    \label{eq:AQ_action_identity}
\end{align}

The coherent photon branch satisfies
\begin{equation}
    \boxed{
    I_{\gamma}
    \equiv
    \frac{E_{\gamma}}{\omega_{\gamma}}
    =
    \hbar.
    }
    \label{eq:AQ_photon_action}
\end{equation}

Second, the present appendix introduces one explicit microscopic extension:
the established transverse Soft-Floor modal coordinate is allowed to couple
weakly to the phase-reference architecture through a
reference-mediated scale-transfer channel.

Third, a dimensionless coefficient
$\mathcal F_{\mathrm R}$ is retained because geometry alone does not determine
the complete carrier-projected noncoherent spectral response.

The derivation therefore has the logical structure
\begin{equation}
    \boxed{
    \text{established TFT microphysics}
    \rightarrow
    \text{reference overlap}
    \rightarrow
    \text{second-order scale kernel}
    \rightarrow
    \Gamma_{\mathrm D}.
    }
    \label{eq:AQ_logic_chain}
\end{equation}


\section{Why Ordinary Amplitude Damping Cannot Produce the Required Drift}

Consider an ordinary linearly damped transverse wave equation
\begin{equation}
    \nabla^2 A
    -
    \frac{1}{c_{\mathrm E}^2}
    \frac{\partial^2 A}{\partial t^2}
    -
    \frac{\gamma_{\mathrm A}}{c_{\mathrm E}^2}
    \frac{\partial A}{\partial t}
    =
    0,
    \label{eq:AQ_linear_damped_wave}
\end{equation}
where $\gamma_{\mathrm A}$ is an ordinary amplitude-damping coefficient.

For the harmonic ansatz
\begin{equation}
    A
    =
    A_0
    e^{i(\mathbf{k}\cdot\mathbf{x}-\omega t)},
\end{equation}
one obtains
\begin{equation}
    \omega^2
    +
    i\gamma_{\mathrm A}\omega
    -
    c_{\mathrm E}^2k^2
    =
    0.
\end{equation}

The positive-frequency root is
\begin{equation}
    \omega
    =
    -\frac{i\gamma_{\mathrm A}}{2}
    +
    \sqrt{
        c_{\mathrm E}^2k^2
        -
        \frac{\gamma_{\mathrm A}^2}{4}
    }.
\end{equation}

For
\begin{equation}
    \gamma_{\mathrm A}
    \ll
    c_{\mathrm E}k,
\end{equation}
this becomes
\begin{equation}
    \omega
    \simeq
    c_{\mathrm E}k
    -
    \frac{i\gamma_{\mathrm A}}{2}
    -
    \frac{\gamma_{\mathrm A}^2}
    {8c_{\mathrm E}k}
    +\cdots.
\end{equation}

The leading imaginary contribution produces amplitude decay. The leading real
frequency remains fixed by the carrier wave number.

Thus
\begin{equation}
    \boxed{
    \text{ordinary fixed-}k\text{ damping}
    \neq
    \text{continuous real-frequency migration}.
    }
    \label{eq:AQ_damping_nogo}
\end{equation}

Consequently, the desired vacuum relaxation rate must not be identified
directly with $\gamma_{\mathrm A}$.

The required reduced photon law is instead
\begin{equation}
    \boxed{
    \frac{d\ln\omega_{\gamma}}{dt}
    =
    -\Gamma_{\mathrm D}.
    }
    \label{eq:AQ_required_frequency_law}
\end{equation}


\section{Self-Similar Photon Scale Coordinate}

Define the dimensionless photon scale coordinate
\begin{equation}
    q
    \equiv
    kr_{\mathrm D}.
    \label{eq:AQ_qdef}
\end{equation}

The coherent self-similar photon family may be represented schematically by
\begin{equation}
    \mathbf{u}_q(\mathbf{x})
    =
    r_{\mathrm D}q\,
    \mathbf{U}
    \left(
        \frac{q\mathbf{x}}{r_{\mathrm D}}
    \right).
    \label{eq:AQ_selfsimilar_photon}
\end{equation}

Introduce
\begin{equation}
    \boldsymbol{\xi}
    =
    \frac{q\mathbf{x}}{r_{\mathrm D}}.
\end{equation}

The associated strain scales as
\begin{equation}
    \boldsymbol{\varepsilon}_q
    =
    q^2
    \boldsymbol{\mathcal E}(\boldsymbol{\xi}),
\end{equation}
whereas the volume element transforms according to
\begin{equation}
    d^3x
    =
    r_{\mathrm D}^{3}
    q^{-3}
    d^3\xi.
\end{equation}

The elastic energy therefore takes the form
\begin{align}
    E_{\gamma}(q)
    &=
    \frac{1}{2}
    \int
    \boldsymbol{\varepsilon}_q
    :
    \mathbf{C}
    :
    \boldsymbol{\varepsilon}_q
    \,d^3x
    \\
    &=
    qE_*,
    \label{eq:AQ_Elinearq}
\end{align}
with
\begin{equation}
    E_*
    =
    \frac{r_{\mathrm D}^{3}}{2}
    \int
    \boldsymbol{\mathcal E}
    :
    \mathbf{C}
    :
    \boldsymbol{\mathcal E}
    \,d^3\xi.
\end{equation}

For the one-action photon family,
\begin{equation}
    E_{\gamma}
    =
    \hbar\omega_{\gamma}.
    \label{eq:AQ_Ehbaromega}
\end{equation}

Define the logarithmic scale coordinate
\begin{equation}
    s
    \equiv
    \ln q.
    \label{eq:AQ_sdef}
\end{equation}

Then
\begin{equation}
    \dot s
    =
    \frac{\dot q}{q}.
\end{equation}

Let
$\boldsymbol{\Phi}_{\gamma}$
denote the complete coherent photon configuration. The tangent to the
self-similar photon family is
\begin{equation}
    \boxed{
    \boldsymbol{\Psi}_s
    \equiv
    \frac{
        \partial\boldsymbol{\Phi}_{\gamma}
    }{
        \partial s
    }
    =
    q
    \frac{
        \partial\boldsymbol{\Phi}_{\gamma}
    }{
        \partial q
    }.
    }
    \label{eq:AQ_scale_tangent}
\end{equation}

The proposed weak nonconservative interaction is required to project primarily
onto this tangent.

It is therefore a scale-transfer process rather than an ordinary attenuation
operator acting orthogonally to the coherent photon family.


\section{Established Role of the FCC Modal Area}

The FCC contact homogenization gives
\begin{equation}
    \boxed{
    A_{\mathrm{ct}}^{(\mathrm{modal})}
    =
    \frac{M_{\mathrm D}}{\mu_{\mathrm D}}
    =
    2\pi\ell_{\mathrm S}^{2}.
    }
    \label{eq:AQ_modal_area}
\end{equation}

The quantity
$A_{\mathrm{ct}}^{(\mathrm{modal})}$
has a specific established role.

The microscopic FCC contact law is expressed in terms of interface
coefficients, while the primitive-cell reduced transverse coordinate is
expressed through integrated stiffness and nodal inertia.

The modal inertia is
\begin{equation}
    m_{\mathrm{ct}}^{(\mathrm{modal})}
    =
    \mu_{\mathrm D}
    A_{\mathrm{ct}}^{(\mathrm{modal})}
    =
    M_{\mathrm D}.
    \label{eq:AQ_modal_mass}
\end{equation}

Similarly,
\begin{equation}
    k_{\mathrm{ct}}^{(\mathrm{modal})}
    =
    \kappa_{\mathrm c}
    A_{\mathrm{ct}}^{(\mathrm{modal})}.
    \label{eq:AQ_modal_stiffness}
\end{equation}

Thus
\begin{equation}
    \boxed{
    A_{\mathrm{ct}}^{(\mathrm{modal})}
    =
    \text{surface-to-cell normalization area of the reduced transverse mode}.
    }
\end{equation}

It is not asserted to be a literal circular contact patch or a geometrically
isolated active fraction of every FCC face.

The present derivation preserves this interpretation.


\section{Reference-Mediated Scale-Transfer Closure}

The Hard Floor remains the zero-inertia phase-reference and structural-record
sector.

It is not assigned:

\begin{itemize}
    \item an independently propagating photon mode,
    \item an additive photon inertia,
    \item or a dissipative energy-absorption role.
\end{itemize}

The minimal new closure is that the same normalized transverse coordinate used
by the Soft-Floor photon may weakly sample the fixed substrate reference.

Define the transverse reference-coupled modal measure by
\begin{equation}
    \boxed{
    A_{\mathrm R}
    \equiv
    A_{\mathrm{ct}}^{(\mathrm{modal})}.
    }
    \label{eq:AQ_AR}
\end{equation}

Equation~\eqref{eq:AQ_AR} does not identify the modal area with a literal
Hard-Floor aperture.

It states only that the reference-mediated coupling is projected through the
same normalized transverse coordinate already established for the photon
carrier.

The additional microscopic closure assigns one substrate correlation depth
$\ell_{\mathrm S}$ to this reference support:
\begin{equation}
    \boxed{
    V_{\mathrm R}
    \equiv
    A_{\mathrm R}\ell_{\mathrm S}.
    }
    \label{eq:AQ_VR}
\end{equation}

Using Equation~\eqref{eq:AQ_modal_area},
\begin{equation}
    \boxed{
    V_{\mathrm R}
    =
    2\pi\ell_{\mathrm S}^{3}.
    }
    \label{eq:AQ_VR_explicit}
\end{equation}

The Soft-Floor node volume is
\begin{equation}
    \boxed{
    V_{\mathrm D}
    =
    \frac{4\pi}{3}
    r_{\mathrm D}^{3}.
    }
    \label{eq:AQ_VD}
\end{equation}

For normalized lowest-order support functions, define the geometric overlap
amplitude
\begin{equation}
    \eta_{\mathrm R}
    \equiv
    \sqrt{
        \frac{V_{\mathrm R}}{V_{\mathrm D}}
    }.
    \label{eq:AQ_eta_def}
\end{equation}

Substitution gives
\begin{equation}
    \boxed{
    \eta_{\mathrm R}
    =
    \sqrt{\frac{3}{2}}
    \left(
        \frac{\ell_{\mathrm S}}{r_{\mathrm D}}
    \right)^{3/2}.
    }
    \label{eq:AQ_eta}
\end{equation}

This quantity is an overlap amplitude between normalized modal supports.

It is not a collision probability. It is also not a thermodynamic state count,
mode degeneracy, projector rank, or heat-capacity fraction of the receiving
sector.


\section{Second-Order Number-Preserving Scale Coupling}

A number-preserving photon transition cannot terminate in the Hard Floor.

The phase-reference sector supplies only the intermediate structural
correlation. The final energy must remain within the active Soft-Floor system.

The lowest-order schematic process is therefore
\begin{equation}
    \begin{aligned}
        \gamma_q
        &\longrightarrow
        \text{reference-correlated intermediate state}\\
        &\longrightarrow
        \gamma_{q-dq}
        +\text{noncoherent active-vacuum degrees of freedom},
    \end{aligned}
    \label{eq:AQ_second_order_process}
\end{equation}
where the noncoherent degrees of freedom receive the energy transferred out of
the coherent photon branch. Their subsequent thermalization is not assumed in
the present derivation.

The projected interaction may be written schematically as
\begin{equation}
    H_{\mathrm{int}}^{(\mathrm R)}
    =
    g_0
    \epsilon_{\mathrm R}
    \hat X_s
    \hat B_{\mathrm{AV}},
    \label{eq:AQ_Hint}
\end{equation}
where:

\begin{itemize}
    \item $\hat X_s$ is the photon operator projected onto the scale tangent
    $\boldsymbol{\Psi}_s$;

    \item $\hat B_{\mathrm{AV}}$ is the work-conjugate operator of the
    noncoherent active-vacuum degrees of freedom;

    \item $g_0$ is the unsuppressed microscopic interaction scale;

    \item $\epsilon_{\mathrm R}$ is the dimensionless
    reference-mediated effective coupling.
\end{itemize}

Under the adopted two-reference-overlap-vertex return closure, the effective
coupling amplitude is
\begin{equation}
    \epsilon_{\mathrm R}
    =
    \eta_{\mathrm R}^{2}.
\end{equation}

Hence
\begin{equation}
    \boxed{
    \epsilon_{\mathrm R}
    =
    \frac{3}{2}
    \left(
        \frac{\ell_{\mathrm S}}{r_{\mathrm D}}
    \right)^3.
    }
    \label{eq:AQ_epsilonR}
\end{equation}

The irreversible scale kernel is second order in this effective interaction
amplitude.

Define
\begin{equation}
    P_{\mathrm R}
    \equiv
    \epsilon_{\mathrm R}^{2}.
    \label{eq:AQ_PRdef}
\end{equation}

Then
\begin{equation}
    \boxed{
    P_{\mathrm R}
    =
    \frac{9}{4}
    \left(
        \frac{\ell_{\mathrm S}}{r_{\mathrm D}}
    \right)^6.
    }
    \label{eq:AQ_PR}
\end{equation}

The sixth power is therefore generated by
\begin{equation}
    \boxed{
    \underbrace{
    \left(
        \frac{\ell_{\mathrm S}}{r_{\mathrm D}}
    \right)^3
    }_{\text{two-vertex amplitude}}
    \quad
    \xrightarrow{\;|\cdot|^2\;}
    \quad
    \underbrace{
    \left(
        \frac{\ell_{\mathrm S}}{r_{\mathrm D}}
    \right)^6
    }_{\text{relaxation-kernel weight}}.
    }
    \label{eq:AQ_sixth_power_origin}
\end{equation}

The quantity $P_{\mathrm R}$ is a dimensionless perturbative kernel weight.

It must not be interpreted as the probability of one rare, macroscopic
frequency-changing collision.


\section{Carrier-Projected Noncoherent Correlation}

The geometric factor in Equation~\eqref{eq:AQ_PR} does not by itself generate
irreversibility.

A nonzero relaxation rate requires a continuum of noncoherent active-vacuum
states onto which the photon scale tangent can transfer energy.

Let the reduced logarithmic-scale dynamics be written schematically as
\begin{equation}
    \dot s(t)
    =
    -
    P_{\mathrm R}
    \int_0^{\infty}
    K_{\mathrm R}(t_{\mathrm{lag}})
    \,dt_{\mathrm{lag}}.
    \label{eq:AQ_kernel_equation}
\end{equation}

The reference microscopic two-leg correlation interval is defined as
\begin{equation}
    \boxed{
    t_{\mathrm{corr}}^{(0)}
    \equiv
    \frac{2d_{\mathrm{nn}}}{c_{\mathrm E}}
    =
    \frac{4r_{\mathrm D}}{c_{\mathrm E}}.
    }
    \label{eq:AQ_tau0}
\end{equation}

This is the shortest geometric interval associated with departure from and
return to the coherent nearest-neighbor transfer sequence.

The exact physical correlation duration need not equal this reference value.

Write the kernel as
\begin{equation}
    K_{\mathrm R}(t_{\mathrm{lag}})
    =
    \frac{1}{
        \left(
            t_{\mathrm{corr}}^{(0)}
        \right)^2
    }
    \widehat K_{\mathrm R}
    \left(
        \frac{t_{\mathrm{lag}}}{
            t_{\mathrm{corr}}^{(0)}
        }
    \right).
    \label{eq:AQ_kernel_scaled}
\end{equation}

Define the dimensionless carrier--noncoherent correlation factor
\begin{equation}
    \boxed{
    \mathcal F_{\mathrm R}
    \equiv
    \int_0^{\infty}
    \widehat K_{\mathrm R}(u)\,du,
    \qquad
    u
    =
    \frac{t_{\mathrm{lag}}}{t_{\mathrm{corr}}^{(0)}}.
    }
    \label{eq:AQ_FRdef}
\end{equation}

The factor $\mathcal F_{\mathrm R}$ contains the normalized microscopic
information not determined by the geometric overlap alone, including:

\begin{itemize}
    \item the intermediate-state energy denominator;

    \item the scale-projected active-vacuum spectral density;

    \item phase and registry participation;

    \item the exact temporal correlation function;

    \item the detailed projection of the interaction onto
    $\boldsymbol{\Psi}_s$;

    \item and the physical correlation duration relative to the reference
    two-leg interval.
\end{itemize}

The carrier-projected spectral density entering
$\mathcal F_{\mathrm R}$ is not the thermodynamic density of states of the
receiving sector and therefore does not determine its state multiplicity,
thermal energy, or heat capacity.

The reduced scale law becomes
\begin{equation}
    \dot s
    =
    -
    \frac{
        P_{\mathrm R}
        \mathcal F_{\mathrm R}
    }{
        t_{\mathrm{corr}}^{(0)}
    }.
\end{equation}

Identifying
\begin{equation}
    \dot s
    =
    -\Gamma_{\mathrm D},
\end{equation}
one obtains
\begin{equation}
    \Gamma_{\mathrm D}
    =
    \frac{
        P_{\mathrm R}
        \mathcal F_{\mathrm R}
    }{
        t_{\mathrm{corr}}^{(0)}
    }.
    \label{eq:AQ_gamma_kernel}
\end{equation}

Substituting Equations~\eqref{eq:AQ_PR} and
\eqref{eq:AQ_tau0} yields
\begin{equation}
    \boxed{
    \Gamma_{\mathrm D}
    =
    \frac{9}{16}
    \mathcal F_{\mathrm R}
    \frac{c_{\mathrm E}}{r_{\mathrm D}}
    \left(
        \frac{\ell_{\mathrm S}}{r_{\mathrm D}}
    \right)^6.
    }
    \label{eq:AQ_gamma_final}
\end{equation}

Equation~\eqref{eq:AQ_gamma_final} is the principal conditional microscopic
closure of this appendix.


\section{Required Continuum and the Conservative Limit}

A finite set of conservative internal modes cannot produce an irreversible
constant relaxation coefficient in the asymptotic long-time limit.

Let
\begin{equation}
    \Sigma_s^{\mathrm R}(\Omega)
\end{equation}
denote the retarded self-energy of the photon scale coordinate after
noncoherent active-vacuum degrees of freedom have been projected out.

A finite Markovian relaxation coefficient requires
\begin{equation}
    \boxed{
    \eta_{s\mathrm R}
    \equiv
    \lim_{\Omega\rightarrow0}
    \frac{
        -\operatorname{Im}
        \Sigma_s^{\mathrm R}(\Omega)
    }{
        \Omega
    }
    >
    0.
    }
    \label{eq:AQ_ohmic_condition}
\end{equation}

A finite pure-point spectrum of conservative oscillators does not satisfy this
condition in the required irreversible sense.

Therefore,
\begin{equation}
    \boxed{
    \text{pure conservative coherent sector}
    \Longrightarrow
    \mathcal F_{\mathrm R}=0
    \Longrightarrow
    \Gamma_{\mathrm D}=0.
    }
    \label{eq:AQ_conservative_limit}
\end{equation}

This reproduces the previously established conservative coherent-vacuum
result. It applies to the completed conservative branch and does not by itself
determine the response of a real nonideal State-I Active Vacuum.

A finite vacuum relaxation rate requires both
\begin{equation}
    P_{\mathrm R}\neq0
\end{equation}
and
\begin{equation}
    \mathcal F_{\mathrm R}\neq0.
\end{equation}

The first measures the microscopic weakness of the coupling.

The second determines whether the active vacuum actually supplies a
carrier-projected irreversible continuum.


\section{Unit-Reference Single-Correlation Closure}

The simplest local normalization is
\begin{equation}
    \boxed{
    \mathcal F_{\mathrm R}=1.
    }
    \label{eq:AQ_FRone}
\end{equation}

This is not claimed to be derived from the complete nonlinear microscopic
spectrum.

It is the unit-reference single-correlation closure.

Under Equation~\eqref{eq:AQ_FRone},
\begin{equation}
    \boxed{
    \Gamma_{\mathrm D}^{(0)}
    =
    \frac{9}{16}
    \frac{c_{\mathrm E}}{r_{\mathrm D}}
    \left(
        \frac{\ell_{\mathrm S}}{r_{\mathrm D}}
    \right)^6.
    }
    \label{eq:AQ_gamma0}
\end{equation}

Using the presently adopted TFT microscopic quantities
\begin{align}
    c_{\mathrm E}
    &=
    2.99792458\times10^8
    \ {\rm m\,s^{-1}},
    \\
    \ell_{\mathrm S}
    &\simeq
    1.616255\times10^{-35}
    \ {\rm m},
    \\
    r_{\mathrm D}
    &\simeq
    7.4504\times10^{-27}
    \ {\rm m},
\end{align}
gives
\begin{equation}
    \frac{\ell_{\mathrm S}}{r_{\mathrm D}}
    \simeq
    2.16935\times10^{-9}.
\end{equation}

The single-reference overlap is
\begin{equation}
    \eta_{\mathrm R}
    \simeq
    1.2375\times10^{-13}.
\end{equation}

The effective two-vertex amplitude is
\begin{equation}
    \epsilon_{\mathrm R}
    \simeq
    1.5314\times10^{-26}.
\end{equation}

The corresponding second-order kernel weight is
\begin{equation}
    \boxed{
    P_{\mathrm R}
    \simeq
    2.3451\times10^{-52}.
    }
    \label{eq:AQ_PR_numeric}
\end{equation}

The reference correlation frequency is
\begin{equation}
    \frac{1}{
        t_{\mathrm{corr}}^{(0)}
    }
    =
    \frac{
        c_{\mathrm E}
    }{
        4r_{\mathrm D}
    }
    \simeq
    1.0060\times10^{34}
    \ {\rm s^{-1}}.
\end{equation}

Their product gives
\begin{equation}
    \boxed{
    \Gamma_{\mathrm D}^{(0)}
    \simeq
    2.3591\times10^{-18}
    \ {\rm s^{-1}}.
    }
    \label{eq:AQ_gamma_numeric}
\end{equation}

No astronomical observable is used in obtaining
Equation~\eqref{eq:AQ_gamma_numeric}. This remains a unit-reference
evaluation: it assumes $\mathcal F_{\mathrm R}=1$ and inherits the adopted
calibration of $r_{\mathrm D}$.


\section{Downstream Calibrated Normalization}

Only after evaluating the unit-reference microscopic rate do we introduce the present low-redshift empirical normalization. Define

\begin{equation}
    \Gamma_{\mathrm D}^{(\mathrm{cal})}
    \equiv
    H_0^{(\mathrm{obs})}.
\end{equation}

Using the value employed in the cosmological sector,

\begin{equation}
    \Gamma_{\mathrm D}^{(\mathrm{cal})}
    \simeq
    2.3671\times10^{-18}\ \mathrm{s}^{-1}.
\end{equation}

The corresponding carrier-response normalization is

\begin{equation}
    F_{\mathrm R}^{(\mathrm{cal})}
    =
    \frac{
        \Gamma_{\mathrm D}^{(\mathrm{cal})}
    }{
        \Gamma_{\mathrm D}^{(0)}
    }.
\end{equation}

Numerically,

\begin{equation}
    F_{\mathrm R}^{(\mathrm{cal})}
    \simeq
    1.00339.
\end{equation}

Thus, the unit-reference rate differs from the present calibrated normalization by approximately $0.34\%$. This downstream normalization leaves the microscopic unit-reference construction unchanged.

\section{Equivalent Fundamental-Constant Representation}

The substrate spacing satisfies
\begin{equation}
    \ell_{\mathrm S}^{2}
    =
    \frac{
        \hbar G
    }{
        c_{\mathrm E}^{3}
    }.
    \label{eq:AQ_lS}
\end{equation}

Hence
\begin{equation}
    \ell_{\mathrm S}^{6}
    =
    \frac{
        \hbar^3G^3
    }{
        c_{\mathrm E}^{9}
    }.
\end{equation}

Equation~\eqref{eq:AQ_gamma_final} may therefore be written
\begin{equation}
    \boxed{
    \Gamma_{\mathrm D}
    =
    \frac{9}{16}
    \mathcal F_{\mathrm R}
    \frac{
        \hbar^3G^3
    }{
        c_{\mathrm E}^{8}
        r_{\mathrm D}^{7}
    }.
    }
    \label{eq:AQ_gamma_fundamental}
\end{equation}

The dimensions are
\begin{equation}
    [\Gamma_{\mathrm D}]
    =
    {\rm s^{-1}}.
\end{equation}

The result exposes the strong scale sensitivity
\begin{equation}
    \boxed{
    \Gamma_{\mathrm D}
    \propto
    r_{\mathrm D}^{-7}.
    }
    \label{eq:AQ_rD_scaling}
\end{equation}

Taking a formal logarithmic variation before imposing correlations among the
microscopic quantities gives
\begin{equation}
    \frac{
        \delta\Gamma_{\mathrm D}
    }{
        \Gamma_{\mathrm D}
    }
    =
    \frac{
        \delta\mathcal F_{\mathrm R}
    }{
        \mathcal F_{\mathrm R}
    }
    +
    6
    \frac{
        \delta\ell_{\mathrm S}
    }{
        \ell_{\mathrm S}
    }
    -
    7
    \frac{
        \delta r_{\mathrm D}
    }{
        r_{\mathrm D}
    }
    +
    \frac{
        \delta c_{\mathrm E}
    }{
        c_{\mathrm E}
    }.
    \label{eq:AQ_sensitivity}
\end{equation}

When Equation~\eqref{eq:AQ_lS} is imposed, the variations of
$\ell_{\mathrm S}$, $G$, and $c_{\mathrm E}$ are not independent. The numerical
rate remains strongly sensitive to the independent determination of
$r_{\mathrm D}$.

The current TFT construction of $r_{\mathrm D}$ inherits an unresolved
microscopic two-dimensional-to-three-dimensional impedance projection.
Consequently,
\begin{equation}
    \boxed{
    \text{cosmology-independent}
    \;\not\Rightarrow\;
    \text{parameter-free}.
    }
\end{equation}


\section{Continuous Scale Drift Rather Than Rare Events}

The extremely small value of $P_{\mathrm R}$ must not be interpreted as a
rare-event probability.

In the intended coherent-kernel closure, each microscopic correlation interval
produces a weak scale-tangent correction. Let the discrete logarithmic scale
evolution be
\begin{equation}
    s_{n+1}
    =
    s_n
    -
    \mathcal F_{\mathrm R}
    P_{\mathrm R}
    +
    \delta s_n,
    \label{eq:AQ_discrete_drift}
\end{equation}
where $\delta s_n$ is the residual fluctuating component.

After
\begin{equation}
    N_{\mathrm{corr}}
    \simeq
    \frac{
        t
    }{
        t_{\mathrm{corr}}^{(0)}
    }
\end{equation}
correlation intervals,
\begin{equation}
    \langle s(t)\rangle
    =
    s_0
    -
    \Gamma_{\mathrm D}t.
    \label{eq:AQ_s_drift}
\end{equation}

Therefore,
\begin{equation}
    \boxed{
    \frac{
        d\langle\ln q\rangle
    }{
        dt
    }
    =
    -\Gamma_{\mathrm D}.
    }
\end{equation}

If the residual increments are sufficiently weakly correlated and satisfy the
additional variance closure
\begin{equation}
    \left\langle
        (\delta s_n)^2
    \right\rangle
    =
    \mathcal C_{\mathrm R}
    P_{\mathrm R}^{2},
    \label{eq:AQ_noise_assumption}
\end{equation}
with $\mathcal C_{\mathrm R}=O(1)$, the logarithmic-scale diffusion coefficient
is
\begin{equation}
    D_s
    =
    \frac{
        \mathcal C_{\mathrm R}
        P_{\mathrm R}^{2}
    }{
        2t_{\mathrm{corr}}^{(0)}
    }.
\end{equation}

Thus, for $\mathcal F_{\mathrm R}>0$,
\begin{equation}
    \boxed{
    \frac{
        D_s
    }{
        \Gamma_{\mathrm D}
    }
    =
    \frac{
        \mathcal C_{\mathrm R}
    }{
        2\mathcal F_{\mathrm R}
    }
    P_{\mathrm R}.
    }
    \label{eq:AQ_diffusion_ratio}
\end{equation}

Because
\begin{equation}
    P_{\mathrm R}
    \sim
    10^{-52},
\end{equation}
this particular fluctuation closure permits systematic scale drift with
negligible stochastic broadening.

The suppression in Equation~\eqref{eq:AQ_diffusion_ratio} is conditional on
Equation~\eqref{eq:AQ_noise_assumption}, the assumed increment correlations,
and the physical state of the receiver. It does not follow from the smallness
of $P_{\mathrm R}$ alone.

A compound-jump description is not adopted in the present coherent-kernel
closure. Any such alternative must be assessed through its independently
derived first and second transfer moments rather than excluded solely by its
statistical classification.

When $D_s$ is not negligible, the drift of $\langle s\rangle$ and the
fractional decay of the mean photon energy need not coincide exactly. Outside
the drift-dominated regime, the physical rate must be defined by the
energy-transfer first moment.


\section{Energy Transfer and Real-Frequency Evolution}

The projected scale-transfer law gives
\begin{equation}
    \boxed{
    \left.\dot E_{\gamma}\right|_{\mathrm{rel}}
    =
    -\Gamma_{\mathrm D}E_{\gamma}.
    }
    \label{eq:AQ_energy_decay}
\end{equation}

After reversible interaction energy and temporary dressing have been accounted
for, the energy removed from the coherent photon branch is transferred to
finite-energy noncoherent degrees of freedom of the Active Vacuum. This
energy-transfer statement does not by itself establish thermalization, a CMB
receiver occupation, or a temperature. The smooth reduced law also does not
imply that a fractional quantum of action is deposited on each individual FCC
node.

For a photon remaining on the same coherent one-action branch,
\begin{equation}
    I_{\gamma}
    =
    \frac{
        E_{\gamma}
    }{
        \omega_{\gamma}
    }
    =
    \hbar.
\end{equation}

The invariant $I_{\gamma}=\hbar$ applies to the surviving coherent photon
branch. The microscopic decomposition of the transferred energy among allowed
noncoherent Active-Vacuum excitations is not determined by the present reduced
scale-transfer law.

Differentiation gives
\begin{equation}
    \dot I_{\gamma}
    =
    \frac{
        \dot E_{\gamma}
    }{
        \omega_{\gamma}
    }
    -
    \frac{
        E_{\gamma}
    }{
        \omega_{\gamma}^{2}
    }
    \dot\omega_{\gamma}.
\end{equation}

Under
\begin{equation}
    \dot I_{\gamma}=0,
\end{equation}
one obtains
\begin{equation}
    \frac{
        \dot E_{\gamma}
    }{
        E_{\gamma}
    }
    =
    \frac{
        \dot\omega_{\gamma}
    }{
        \omega_{\gamma}
    }.
\end{equation}

Hence
\begin{equation}
    \boxed{
    \left.\dot\omega_{\gamma}\right|_{\mathrm{rel}}
    =
    -\Gamma_{\mathrm D}\omega_{\gamma}.
    }
    \label{eq:AQ_omega_evolution}
\end{equation}

For constant $\Gamma_{\mathrm D}$, integration gives
\begin{equation}
    \boxed{
    \omega_{\gamma}(t)
    =
    \omega_{\gamma,0}
    e^{-\Gamma_{\mathrm D}t}.
    }
    \label{eq:AQ_omega_solution}
\end{equation}

On a homogeneous stationary linear branch, the wavelength evolves
correspondingly as
\begin{equation}
    \boxed{
    \lambda_{\gamma}(t)
    =
    \lambda_{\gamma,0}
    e^{+\Gamma_{\mathrm D}t}.
    }
    \label{eq:AQ_lambda_solution}
\end{equation}

This is a real-frequency evolution through the coherent eigenbranch. It is not
obtained from a decaying field amplitude at fixed $k$.


\section{Wave-Vector Transport}

The rate $\Gamma_{\mathrm D}$ and the conservative ray velocity are expressed
using the same evolution parameter $t$. In the homogeneous gravity-neutral
calculation below, this is the inertial time of the resting Active-Vacuum
reference. A rate specified per a different clock must be converted before
combining it with the ray speed.

Let the conservative carrier branch be
\begin{equation}
    \omega_{\gamma}
    =
    \omega_{\mathrm c}(\mathbf{k}).
\end{equation}

For a locally homogeneous stationary segment of the Active Vacuum,
\begin{equation}
    \left.\dot\omega_{\gamma}\right|_{\mathrm{rel}}
    =
    \mathbf{v}_{\mathrm g}
    \cdot
    \left.\dot{\mathbf{k}}\right|_{\mathrm{rel}},
\end{equation}
where
\begin{equation}
    \mathbf{v}_{\mathrm g}
    =
    \nabla_{\mathbf{k}}
    \omega_{\mathrm c}.
\end{equation}

The intended leading scale-transfer channel is direction preserving:
\begin{equation}
    \boxed{
    \left.
    \dot{\hat{\mathbf{k}}}
    \right|_{\mathrm{rel}}
    =
    0.
    }
    \label{eq:AQ_direction_preservation}
\end{equation}

This is a defining condition on the selected mean relaxation channel. Isotropy
alone does not exclude stochastic angular diffusion, whose rate must be
evaluated separately.

Write
\begin{equation}
    \left.\dot{\mathbf{k}}\right|_{\mathrm{rel}}
    =
    \left.\dot k\right|_{\mathrm{rel}}
    \hat{\mathbf{k}}.
\end{equation}

Define the longitudinal group-speed projection
\begin{equation}
    v_{\mathrm{g},\parallel}
    \equiv
    \mathbf{v}_{\mathrm g}
    \cdot
    \hat{\mathbf{k}}.
\end{equation}

Using Equation~\eqref{eq:AQ_omega_evolution},
\begin{equation}
    v_{\mathrm{g},\parallel}
    \left.\dot k\right|_{\mathrm{rel}}
    =
    -\Gamma_{\mathrm D}
    \omega_{\mathrm c}.
\end{equation}

Where $v_{\mathrm{g},\parallel}\neq0$,
\begin{equation}
    \boxed{
    \left.\dot{\mathbf{k}}\right|_{\mathrm{rel}}
    =
    -
    \Gamma_{\mathrm D}
    \frac{
        \omega_{\mathrm c}
    }{
        v_{\mathrm{g},\parallel}
    }
    \hat{\mathbf{k}}.
    }
    \label{eq:AQ_k_general}
\end{equation}

For the homogeneous gravity-neutral long-wave linear branch,
\begin{equation}
    \omega_{\mathrm c}
    \simeq
    c_{\mathrm E}k,
    \qquad
    v_{\mathrm{g},\parallel}
    \simeq
    c_{\mathrm E},
\end{equation}
so
\begin{equation}
    \boxed{
    \left.\dot{\mathbf{k}}\right|_{\mathrm{rel}}
    =
    -\Gamma_{\mathrm D}\mathbf{k}.
    }
    \label{eq:AQ_k_simple}
\end{equation}

Let $f_{\gamma}(\mathbf x,\mathbf k,t)\geq0$ denote photon number per
$\mathrm d^3x\,\mathrm d^3k$ in the reduced long-wave description, summed over
polarizations. Any constant phase-space normalization factor is absorbed
into $f_{\gamma}$. Work on $k=|\mathbf k|>0$, with
$\hat{\mathbf k}=\mathbf k/k$, and with support within the retained long-wave
sector. For a constant nonnegative rate, the corresponding reduced
one-photon kinetic operator is
\begin{equation}
    \boxed{
    \mathcal L_{\mathrm D}f_{\gamma}
    =
    \Gamma_{\mathrm D}
    \nabla_{\mathbf{k}}
    \cdot
    \left(
        \mathbf{k}f_{\gamma}
    \right).
    }
    \label{eq:AQ_LD}
\end{equation}

Together with the leading spatial ray law, the population continuity
equation is
\[
\begin{gathered}
    \dot{\mathbf{x}}=c_{\mathrm E}\hat{\mathbf{k}},
    \qquad
    \dot{\mathbf{k}}=-\Gamma_{\mathrm D}\mathbf{k},\\
    \partial_t f_{\gamma}
    +\nabla_{\mathbf{x}}\cdot(\dot{\mathbf{x}}f_{\gamma})
    +\nabla_{\mathbf{k}}\cdot(\dot{\mathbf{k}}f_{\gamma})=0.
\end{gathered}
\]
Equivalently,
$\partial_t f_{\gamma}+c_{\mathrm E}\hat{\mathbf k}\cdot
\nabla_{\mathbf x}f_{\gamma}=\mathcal L_{\mathrm D}f_{\gamma}$.
This is the deterministic mean-drift specialization; absorption, diffusion,
and returned excitation require their own transport terms.

Equations~\eqref{eq:AQ_k_general} and~\eqref{eq:AQ_k_simple} specify the
required reduced wave-vector transport. A complete microscopic
photon--receiver interaction must reproduce this transport together with the
correct spatial ray equation and mechanical recoil.


\section{Photon Number, Polarization, and Angular Coherence}

For an integrable distribution, define the total retained photon number by
\[
    N_{\mathrm{ph}}(t)
    =\int f_{\gamma}(\mathbf x,\mathbf k,t)
       \,\mathrm d^3x\,\mathrm d^3k.
\]
Assume sufficient regularity and vanishing spatial and wave-vector boundary
fluxes, or periodic spatial boundaries. Integration of the continuity
equation, using Equation~\eqref{eq:AQ_LD}, then gives
\begin{equation}
    \boxed{
    \left.\dot N_{\mathrm{ph}}\right|_{\mathrm{rel}}
    =
    0.
    }
    \label{eq:AQ_number}
\end{equation}

The notation $N_{\mathrm{ph}}$ denotes photon number and is distinct from the
information count $N_{\gamma}$ defined elsewhere in TFT. Photon absorption,
emission, and thermal number-changing processes are separate channels and are
not excluded by Equation~\eqref{eq:AQ_number}.

This conservation statement concerns the whole retained population. The
number within a finite observing aperture also has its spatial boundary
flux; photon-number conservation does not imply that every aperture
receives a fixed count.

For constant $\Gamma_{\mathrm D}$, the characteristics in the homogeneous
linear-carrier limit are
\[
\begin{gathered}
    \mathbf{k}(t)=e^{-\Gamma_{\mathrm D}t}\mathbf{k}_0,
    \qquad
    \hat{\mathbf{k}}(t)=\hat{\mathbf{k}}_0,\\
    \mathbf{x}(t)=\mathbf{x}_0+c_{\mathrm E}t\hat{\mathbf{k}}_0.
\end{gathered}
\]
The spatial diagonal block of the characteristic derivative is the identity,
and its wave-vector diagonal block is
$e^{-\Gamma_{\mathrm D}t}\mathbf I$. The derivative is block triangular, so
\[
    \det\!\left[
        \frac{\partial(\mathbf x,\mathbf k)}
             {\partial(\mathbf x_0,\mathbf k_0)}
    \right]
    =e^{-3\Gamma_{\mathrm D}t},
    \qquad
    \nabla_{\mathbf x}\cdot\dot{\mathbf x}
    +\nabla_{\mathbf k}\cdot\dot{\mathbf k}
    =-3\Gamma_{\mathrm D}.
\]
The off-diagonal dependence of $\mathbf x(t)$ on $\mathbf k_0$ does not
alter this determinant. Along each characteristic,
$\mathrm d f_{\gamma}/\mathrm dt=3\Gamma_{\mathrm D}f_{\gamma}$, yielding
\[
    f_{\gamma}(\mathbf x,\mathbf k,t)
    =e^{3\Gamma_{\mathrm D}t}
      f_{\gamma,0}\!\left(
          \mathbf x-c_{\mathrm E}t\hat{\mathbf k},
          e^{\Gamma_{\mathrm D}t}\mathbf k
      \right),
\]
where $f_{\gamma,0}$ is the initial distribution. The density increase
compensates exactly for contraction of the reduced phase-space volume.
For spatially uniform time-dependent $\Gamma_{\mathrm D}(t)$, replace
$\Gamma_{\mathrm D}t$ in these scaling factors by
$\int_0^t\Gamma_{\mathrm D}(t')\,\mathrm dt'$. Spatially or spectrally
dependent rates require their own characteristic solution.

For $\Gamma_{\mathrm D}>0$, this reduced flow is not a canonical Hamiltonian
flow of an isolated photon subsystem in $(\mathbf x,\hbar\mathbf k)$.
A complete canonical photon--receiver system can preserve its full
phase-space volume while a reduced photon description contracts after
receiving variables are eliminated. The present drift operator does not
supply that microscopic completion: its projection must also reproduce the
spatial propagation, receiver recoil, and fluctuations. Number conservation
alone therefore does not impose collisionless phase-space transport; the
photometric consequences must retain the actual energy and angular mappings.

The vacuum is assumed parity even and unpolarized at the retained order.
Consequently, the leading mean scale-projected response is required to act
identically on the two transverse polarization states.

In polarization-space notation,
\begin{equation}
    \boxed{
    \Sigma_{ab}^{\mathrm R}
    =
    \Sigma_{\mathrm R}^{\mathrm R}
    \delta_{ab}
    }
    \label{eq:AQ_pol_selfenergy}
\end{equation}
at leading isotropic order.

Hence the differential mean polarization relaxation satisfies
\begin{equation}
    \boxed{
    \Delta\Gamma_{\mathrm{pol}}
    =
    0
    }
    \label{eq:AQ_pol_zero}
\end{equation}
at this order.

This does not by itself exclude stochastic polarization conversion.

The leading selection rule is schematically
\begin{equation}
    \boxed{
    \gamma(
        \mathbf{k},
        \lambda_{\mathrm{pol}}
    )
    \rightarrow
    \gamma(
        \mathbf{k}-d\mathbf{k},
        \lambda_{\mathrm{pol}}
    )
    +
    \mathrm{AV},
    }
    \label{eq:AQ_selection}
\end{equation}
where $\mathrm{AV}$ denotes the noncoherent Active-Vacuum degrees of freedom
receiving the transferred coherent energy, and
\begin{equation}
    d\mathbf{k}
    \parallel
    \mathbf{k}.
\end{equation}

For the selected coherent channel to dominate, the orthogonal rates must
satisfy
\begin{equation}
    \boxed{
    \Gamma_{\mathrm{angle}},
    \Gamma_{\mathrm{pol}},
    \Gamma_{\mathrm{destroy}},
    \Gamma_{\mathrm{secondary}}
    \ll
    \Gamma_{\mathrm D}.
    }
    \label{eq:AQ_suppression_requirements}
\end{equation}

These are direct falsifiability requirements of the proposed microscopic
channel. They are not established solely by photon-number preservation or
isotropy.


\section{Energy and Momentum Ledger}

The Hard Floor is not the dissipative receiving sector. It supplies the fixed
reference scale and structural phase architecture.

After reversible dressing and temporary interaction energy have been included
once and coarse-grained consistently, the secular local energy ledger is
\begin{equation}
    \boxed{
    \left.\dot E_{\gamma}\right|_{\mathrm{rel}}
    +
    \left.\dot E_{\mathrm{AV}}\right|_{\mathrm{rec}}
    =
    0.
    }
    \label{eq:AQ_energy_conservation}
\end{equation}

Using
\begin{equation}
    \left.\dot E_{\gamma}\right|_{\mathrm{rel}}
    =
    -\Gamma_{\mathrm D}E_{\gamma},
\end{equation}
one obtains
\begin{equation}
    \boxed{
    \left.\dot E_{\mathrm{AV}}\right|_{\mathrm{rec}}
    =
    +\Gamma_{\mathrm D}E_{\gamma}.
    }
    \label{eq:AQ_AV_energy}
\end{equation}

Here $E_{\mathrm{AV}}$ denotes only the participating finite-energy
Active-Vacuum receiving sector. It is not the total energy of the Structural
Reservoir and is not identified here with thermal-equilibrium energy.

The Hard-Floor energy deposition is
\begin{equation}
    \boxed{
    \left.\dot E_{\mathrm{HF}}\right|_{\mathrm{rel}}
    =
    0.
    }
    \label{eq:AQ_HF_energy}
\end{equation}

For the long-wave photon,
\begin{equation}
    \mathbf{p}_{\gamma}
    =
    \hbar\mathbf{k}.
\end{equation}

Equation~\eqref{eq:AQ_k_simple} gives
\begin{equation}
    \left.\dot{\mathbf{p}}_{\gamma}\right|_{\mathrm{rel}}
    =
    -\Gamma_{\mathrm D}\mathbf{p}_{\gamma}.
\end{equation}

The corresponding reduced quasimomentum requirement is
\begin{equation}
    \boxed{
    \left.\dot{\mathbf P}_{\mathrm{AV}}\right|_{\mathrm{rec}}
    =
    +\Gamma_{\mathrm D}
    \mathbf{p}_{\gamma}.
    }
    \label{eq:AQ_momentum_balance}
\end{equation}

A complete microscopic derivation must relate this reduced quasimomentum
transfer to the physical FCC--Condensate stress and mechanical recoil.

For the same homogeneous gravity-neutral population used above, assume
finite energy and wave-momentum moments and vanishing moment boundary
fluxes. Define
\[
\begin{aligned}
    \mathcal E_{\gamma}(t)
    &=\int\hbar c_{\mathrm E}k\,f_{\gamma}
        \,\mathrm d^3x\,\mathrm d^3k,\\
    \mathbf P_{\mathrm w}(t)
    &=\int\hbar\mathbf k\,f_{\gamma}
        \,\mathrm d^3x\,\mathrm d^3k.
\end{aligned}
\]
Here $\mathcal E_{\gamma}$ is the total energy of the retained photon
population, distinct from the single-photon $E_{\gamma}$, and
$\mathbf P_{\mathrm w}$ is its canonical wave-quasimomentum moment.
Integration by parts in the continuity equation gives
\[
\begin{aligned}
    \frac{\mathrm d\mathcal E_{\gamma}}{\mathrm dt}
    &=-\Gamma_{\mathrm D}
      \int f_{\gamma}\,
      \mathbf k\cdot\nabla_{\mathbf k}(\hbar c_{\mathrm E}k)
      \,\mathrm d^3x\,\mathrm d^3k
      =-\Gamma_{\mathrm D}\mathcal E_{\gamma},\\
    \frac{\mathrm d\mathbf P_{\mathrm w}}{\mathrm dt}
    &=-\Gamma_{\mathrm D}\mathbf P_{\mathrm w}.
\end{aligned}
\]
Thus the specified drift preserves photon number while reducing the
population energy and wave-quasimomentum moment. The energy decrement
must enter the Active-Vacuum receiving ledger. The wave-quasimomentum
moment is not automatically the complete mechanical momentum of the
scaffold, condensate, and dressed carrier; its conversion to mechanical
recoil remains the same requirement stated above. These moment equations
are consequences of the retained drift, not an independent calculation of
$F_{\mathrm R}$ or of a physical nonzero $\Gamma_{\mathrm D}$.

\section{Leading Achromaticity}

At fixed local background conditions,
Equation~\eqref{eq:AQ_gamma_final} may be written as
\begin{equation}
    \boxed{
    \Gamma_{\mathrm D}(q)
    =
    \Gamma_{\mathrm D}^{(0)}
    \mathcal F_{\mathrm R}(q)
    =
    \frac{
        P_{\mathrm R}
    }{
        t_{\mathrm{corr}}^{(0)}
    }
    \mathcal F_{\mathrm R}(q).
    }
    \label{eq:AQ_rate_factorization}
\end{equation}

Here
\begin{equation}
    \Gamma_{\mathrm D}^{(0)}
    =
    \frac{
        P_{\mathrm R}
    }{
        t_{\mathrm{corr}}^{(0)}
    }
\end{equation}
is the unit-reference rate defined by
\(\mathcal F_{\mathrm R}=1\).
It is distinct from the physical infrared value of the relaxation rate.

Under fixed
\(c_{\mathrm E}\),
\(r_{\mathrm D}\),
and
\(\ell_{\mathrm S}\),
all carrier-scale dependence is contained in
\(\mathcal F_{\mathrm R}(q)\).
Any dependence arising from variation of the local background quantities
themselves must be retained separately and is not absorbed implicitly into
the present fixed-background factorization.

The geometric prefactor in
Equation~\eqref{eq:AQ_rate_factorization}
contains no explicit external photon scale \(q\).
If the carrier-projected receiving response approaches a finite,
nonzero, approximately \(q\)-independent infrared plateau,
\begin{equation}
    \mathcal F_{\mathrm R}(q)
    \longrightarrow
    \mathcal F_{\mathrm R}^{(\mathrm{IR})}
    \qquad
    \text{as}
    \qquad
    q\longrightarrow0,
\end{equation}
define the corresponding physical infrared rate by
\begin{equation}
    \Gamma_{\mathrm D}^{(\mathrm{IR})}
    \equiv
    \Gamma_{\mathrm D}^{(0)}
    \mathcal F_{\mathrm R}^{(\mathrm{IR})}.
\end{equation}

The leading long-wave response is then
\begin{equation}
    \boxed{
    \Gamma_{\mathrm D}(q)
    \simeq
    \Gamma_{\mathrm D}^{(\mathrm{IR})}.
    }
    \label{eq:AQ_achromatic_leading}
\end{equation}

Finite-wave-number, registry, receiver-bandwidth, or interaction-projection
corrections may instead produce
\begin{equation}
    \Gamma_{\mathrm D}(q)
    =
    \Gamma_{\mathrm D}^{(\mathrm{IR})}
    \left[
        1
        +
        O(q^m)
    \right],
    \qquad
    m>0.
    \label{eq:AQ_chromatic_correction}
\end{equation}

The exponent \(m\) and its coefficient must be obtained from the physical
receiving response. Geometry alone does not prove either a nonzero infrared
plateau or leading achromaticity.

Within the constant-response plateau, the notation
\(\Gamma_{\mathrm D}\) may be used for the constant physical value
\(\Gamma_{\mathrm D}^{(\mathrm{IR})}\). The leading fractional evolution is
then
\begin{equation}
    \boxed{
    \frac{1}{\omega_{\gamma}}
    \left.
    \frac{
        d\omega_{\gamma}
    }{
        dt
    }
    \right|_{\mathrm{rel}}
    =
    -\Gamma_{\mathrm D}.
    }
    \label{eq:AQ_fractional_achromaticity}
\end{equation}

Thus, the fractional coherent-frequency evolution is independent of photon
frequency at the retained order only when the carrier-projected receiving
response possesses the required approximately constant infrared plateau.

\section{Effective Loss-Tangent Representation}

For weak coherent energy transfer, an effective quality factor may be defined
through
\begin{equation}
    Q_{\mathrm{eff}}^{-1}
    \equiv
    \frac{
        \Gamma_{\mathrm D}
    }{
        \omega
    }.
\end{equation}

An equivalent reduced loss-tangent notation is
\begin{equation}
    \boxed{
    \tan\delta_{\mathrm{eff}}(\omega)
    \simeq
    \frac{
        \Gamma_{\mathrm D}
    }{
        \omega
    }.
    }
    \label{eq:AQ_loss_tangent}
\end{equation}

This expression is not interpreted as the local material loss tangent of an
ordinary Kelvin--Voigt or Maxwell element acting directly on photon amplitude.

The fundamental nonconservative object is the scale-projected retarded kernel.
Causality requires
\begin{equation}
    K_{\mathrm R}(t_{\mathrm{lag}})
    =
    0,
    \qquad
    t_{\mathrm{lag}}<0.
\end{equation}

The full frequency-dependent response must obey the corresponding
Kramers--Kronig relations.

Equation~\eqref{eq:AQ_loss_tangent} is therefore a reduced weak-loss
representation over the photon-frequency domain in which the fractional
scale-relaxation rate is approximately constant.


\section{Relaxation-Induced Frequency-Shift Component}

The microscopic rate developed in this appendix defines the
photon-propagation relaxation component of a frequency shift.

For a locally quasistatic ray segment,
\begin{equation}
    \frac{
        d\ln\omega_{\gamma}
    }{
        dt
    }
    =
    -\Gamma_{\mathrm D}.
    \label{eq:AQ_local_relaxation}
\end{equation}

Along a ray,
\begin{equation}
    dt
    =
    \frac{
        d\ell
    }{
        v_{\mathrm{ray}}
    }.
\end{equation}

Therefore,
\begin{equation}
    d\ln\omega_{\gamma}
    =
    -
    \frac{
        \Gamma_{\mathrm D}(\ell)
    }{
        v_{\mathrm{ray}}(\ell)
    }
    d\ell.
\end{equation}

Define the relaxation-induced frequency-shift factor by
\begin{equation}
    \boxed{
    1+z_{\mathrm D}
    \equiv
    \frac{
        \omega_{\mathrm{in}}
    }{
        \omega_{\mathrm{out}}
    }.
    }
    \label{eq:AQ_zD_def}
\end{equation}

Integration gives
\begin{equation}
    \boxed{
    \ln(1+z_{\mathrm D})
    =
    \int_{\mathrm{path}}
    \frac{
        \Gamma_{\mathrm D}(\ell)
    }{
        v_{\mathrm{ray}}(\ell)
    }
    \,d\ell.
    }
    \label{eq:AQ_zD_general}
\end{equation}

For constant $\Gamma_{\mathrm D}$ and
$v_{\mathrm{ray}}\simeq c_{\mathrm E}$,
\begin{equation}
    \boxed{
    1+z_{\mathrm D}
    =
    \exp
    \left(
        \frac{
            \Gamma_{\mathrm D}r
        }{
            c_{\mathrm E}
        }
    \right).
    }
    \label{eq:AQ_zD_constant}
\end{equation}

At small relaxation depth,
\begin{equation}
    \frac{
        \Gamma_{\mathrm D}r
    }{
        c_{\mathrm E}
    }
    \ll1,
\end{equation}
this reduces to
\begin{equation}
    \boxed{
    z_{\mathrm D}
    \simeq
    \frac{
        \Gamma_{\mathrm D}r
    }{
        c_{\mathrm E}
    }.
    }
    \label{eq:AQ_zD_low}
\end{equation}

The symbol $z_{\mathrm D}$ is deliberately distinguished from the total
astronomically observed redshift $z_{\mathrm{obs}}$.


\section{Separation of Relaxation Redshift from Endpoint Contributions}

The TFT universe is not assumed to be locally motionless.

Global Structural Reservoir staticity does not eliminate local motion of
matter, localized defects, galaxies, compact objects, or FCC--SC
phase-boundary structures.

A moving emitter and observer can therefore generate a kinematic frequency
shift independently of the propagation relaxation derived above. Different
gravitational loading at the emitter and observer can also produce an endpoint
clock-frequency difference.

Define the corresponding factors by
\begin{equation}
    1+z_{\mathrm{kin}},
    \qquad
    1+z_{\mathrm{grav}}.
\end{equation}

When the endpoint transformations and the propagation relaxation are
consistently factorized, their frequency ratios multiply:
\begin{equation}
    \boxed{
    1+z_{\mathrm{obs}}
    =
    (1+z_{\mathrm{kin}})
    (1+z_{\mathrm{grav}})
    (1+z_{\mathrm D}).
    }
    \label{eq:AQ_total_redshift}
\end{equation}

Equivalently,
\begin{equation}
    \boxed{
    \ln(1+z_{\mathrm{obs}})
    =
    \ln(1+z_{\mathrm{kin}})
    +
    \ln(1+z_{\mathrm{grav}})
    +
    \ln(1+z_{\mathrm D}).
    }
    \label{eq:AQ_total_redshift_log}
\end{equation}

Equation~\eqref{eq:AQ_total_redshift} is an interface relation, not a complete
cosmological closure.

The present appendix derives only the conditional propagation response
$z_{\mathrm D}$. It does not determine $z_{\mathrm{kin}}$,
$z_{\mathrm{grav}}$, the macroscopic source velocity field, the gravitational
loading of the actual emitting region, or the fraction of an observed
astronomical redshift assigned to each contribution.

The full observational decomposition and its source-event timing conditions
are developed in Appendix~\ref{app:time_dilation}. Strongly time-dependent
backgrounds or moving refractive media may require additional path-dependent
frequency terms beyond the stationary factorization above.


\section{The Photon Energy-Relaxation Length and the Structural Reservoir Radius}

For a positive constant physical relaxation rate, the characteristic photon
scale-relaxation length is the TFT Photon Energy-Relaxation Length,
\begin{equation}
    \boxed{
    R_{\mathrm U}
    \equiv
    \frac{
        c_{\mathrm E}
    }{
        \Gamma_{\mathrm D}
    }.
    }
    \label{eq:AQ_RU}
\end{equation}

For constant $\Gamma_{\mathrm D}$,
Equation~\eqref{eq:AQ_zD_constant} becomes
\begin{equation}
    1+z_{\mathrm D}
    =
    e^{r/R_{\mathrm U}}.
\end{equation}

At
\begin{equation}
    r=R_{\mathrm U},
\end{equation}
the coherent photon frequency and energy become
\begin{equation}
    \omega_{\gamma}
    =
    \frac{
        \omega_{\gamma,0}
    }{
        e
    },
\end{equation}
and
\begin{equation}
    E_{\gamma}
    =
    \frac{
        E_{\gamma,0}
    }{
        e
    }.
\end{equation}

Therefore,
\begin{equation}
    \boxed{
    R_{\mathrm U}
    =
    \text{characteristic one-e-fold coherent photon scale-relaxation length}.
    }
    \label{eq:AQ_RU_definition}
\end{equation}

The Photon Energy-Relaxation Length is not:

\begin{itemize}
    \item the physical radius of the TFT universe;

    \item the radius of the Structural Reservoir;

    \item an absolute causal boundary of the substrate;

    \item a photon stopping radius;

    \item a complete photon-energy-removal distance;

    \item a thermalization boundary;

    \item or a requirement that no matter or radiation exists beyond
    $R_{\mathrm U}$.
\end{itemize}

The full structural extent of the TFT universe is described independently by
the Structural Reservoir radius $R_{\mathrm{Res}}$, obtained from the global
phase and mass-conservation sector.

Thus,
\begin{equation}
    \boxed{
    R_{\mathrm U}
    \neq
    R_{\mathrm{Res}}.
    }
    \label{eq:AQ_RU_not_RRes}
\end{equation}

The two quantities describe physically distinct scales:
$R_{\mathrm U}$ characterizes one e-fold of coherent photon scale relaxation,
whereas $R_{\mathrm{Res}}$ characterizes the macroscopic structural extent of
the TFT universe.


\section{Local Dynamics and Constitutive Dependence}

The rate $\Gamma_{\mathrm D}$ is defined in the local rest frame of the active
medium.

{\emergencystretch=2em
Local motion of a source does not modify the microscopic overlap calculation
in Equation~\eqref{eq:AQ_gamma_final} unless that motion changes one of the
local quantities entering\par}
\begin{equation}
    r_{\mathrm D},
    \qquad
    \ell_{\mathrm S},
    \qquad
    \mathcal F_{\mathrm R},
    \qquad
    c_{\mathrm E}.
\end{equation}

Weak gravitational loading alone does not determine
$\mathcal F_{\mathrm R}$. In the completed conservative coherent sector,
\begin{equation}
    \mathcal F_{\mathrm R}=0,
    \qquad
    \Gamma_{\mathrm D}=0.
\end{equation}

For a real nonideal State-I Active Vacuum possessing an independently active
carrier-projected receiving continuum,
\begin{equation}
    \Gamma_{\mathrm D}
    =
    \Gamma_{\mathrm D}^{(0)}
    \mathcal F_{\mathrm R}.
\end{equation}

The specialization
\begin{equation}
    \Gamma_{\mathrm D}
    =
    \Gamma_{\mathrm D}^{(0)}
\end{equation}
holds only under the unit-reference closure
$\mathcal F_{\mathrm R}=1$. It is not the definition of a pristine or merely
weakly loaded vacuum.

Whether large-scale structural evolution produces a systematic dependence of
$\mathcal F_{\mathrm R}$, $r_{\mathrm D}$, or another local constitutive
quantity is not determined here. Such a dependence must be derived from TFT
microphysics rather than fitted retrospectively to astronomical distance data.


\section{What the Present Appendix Establishes}

Conditional on the explicit reference-mediated coupling closure, the present
appendix establishes the following results:

\begin{enumerate}
    \item {\emergencystretch=2em
    Ordinary linear amplitude damping is not the mechanism responsible
    for real photon-frequency migration.\par}

    \item The relevant weak nonconservative coordinate is the tangent to the
    self-similar coherent photon family.

    \item The established transverse modal normalization is
    \begin{equation}
        A_{\mathrm{ct}}^{(\mathrm{modal})}
        =
        2\pi\ell_{\mathrm S}^{2}.
    \end{equation}

    \item Assigning one substrate-reference correlation depth gives
    \begin{equation}
        V_{\mathrm R}
        =
        2\pi\ell_{\mathrm S}^{3}.
    \end{equation}

    \item The normalized single-reference overlap is
    \begin{equation}
        \eta_{\mathrm R}
        =
        \sqrt{\frac32}
        \left(
            \frac{\ell_{\mathrm S}}{r_{\mathrm D}}
        \right)^{3/2}.
    \end{equation}

    \item Under the adopted two-reference-overlap-vertex return closure,
    \begin{equation}
        P_{\mathrm R}
        =
        \frac94
        \left(
            \frac{\ell_{\mathrm S}}{r_{\mathrm D}}
        \right)^6.
    \end{equation}

    \item The resulting conditional rate is
    \begin{equation}
        \boxed{
        \Gamma_{\mathrm D}
        =
        \frac{9}{16}
        \mathcal F_{\mathrm R}
        \frac{
            c_{\mathrm E}
        }{
            r_{\mathrm D}
        }
        \left(
            \frac{
                \ell_{\mathrm S}
            }{
                r_{\mathrm D}
            }
        \right)^6.
        }
    \end{equation}

    \item Under coherent modal-action preservation,
    \begin{equation}
        \frac{
            \dot E_{\gamma}
        }{
            E_{\gamma}
        }
        =
        \frac{
            \dot\omega_{\gamma}
        }{
            \omega_{\gamma}
        }
        =
        -\Gamma_{\mathrm D}.
    \end{equation}

    \item Within the selected leading closure, the mean scale transport is
    constructed to preserve photon number, propagation direction, and
    polarization degeneracy. The suppression of orthogonal channels remains
    an independent falsification requirement.

    \item The microscopic rate defines the propagation-relaxation component
    $z_{\mathrm D}$ rather than automatically defining the total astronomical
    redshift.

    \item For a positive constant physical rate, the Photon Energy-Relaxation Length
    \begin{equation}
        R_{\mathrm U}
        =
        \frac{
            c_{\mathrm E}
        }{
            \Gamma_{\mathrm D}
        }
    \end{equation}
    is the characteristic one-e-fold coherent photon scale-relaxation length
    and is distinct from the Structural Reservoir radius
    $R_{\mathrm{Res}}$.
\end{enumerate}


\section{What the Present Appendix Does Not Establish}

The present derivation does not establish the following stronger statements:

\begin{enumerate}
    \item It does not derive
    \begin{equation}
        \mathcal F_{\mathrm R}=1
    \end{equation}
    from the complete nonlinear FCC--Condensate spectrum.

    \item It does not prove that
    $A_{\mathrm{ct}}^{(\mathrm{modal})}$
    is a literal microscopic contact patch.

    \item It does not yet derive the complete noncoherent Active-Vacuum
    spectral density from the rotational, internal-phase, Condensate, and
    Tension-Column sectors.

    \item It does not derive a complete spatial photon--receiver vertex that
    simultaneously produces real-frequency transfer, the correct spatial ray
    equation, and mechanical FCC--Condensate recoil.

    \item It does not establish the complete higher-order suppression of
    spectral diffusion, angular diffusion, polarization conversion, photon
    destruction, or secondary radiation.

    \item It does not derive that
    $\mathcal F_{\mathrm R}$ is independent of photon scale, position, or the
    physical state of the Active Vacuum.

    \item It does not derive the local galaxy or compact-object velocity field.

    \item It does not assert
    \begin{equation}
        z_{\mathrm{obs}}
        =
        z_{\mathrm D}.
    \end{equation}

    \item It does not derive the gravitational loading of a general emitting
    region or observer.

    \item It does not derive supernova source-envelope time dilation.

    \item It does not derive the complete luminosity-distance relation.

    \item It does not derive BAO geometry.

    \item It does not derive thermalization of the transferred coherent
    photon energy, CMB thermal equilibrium, CMB spectral preservation, or the
    CMB temperature.

    \item It does not identify the Photon Energy-Relaxation Length
    \begin{equation}
        R_{\mathrm U}
    \end{equation}
    with the Structural Reservoir radius
    \begin{equation}
        R_{\mathrm{Res}}.
    \end{equation}

    \item It does not make the numerical value of
    $\Gamma_{\mathrm D}$
    fully parameter free while the independent determination of
    $r_{\mathrm D}$ remains projection dependent.
\end{enumerate}


\section{Falsification Conditions}

The proposed microscopic closure is falsifiable.

First, if a completed FCC--Condensate calculation yields
\begin{equation}
    \lim_{\Omega\rightarrow0}
    \frac{
        -\operatorname{Im}
        \Sigma_s^{\mathrm R}(\Omega)
    }{
        \Omega
    }
    =
    0,
\end{equation}
then the proposed constant Ohmic realization has
\begin{equation}
    \mathcal F_{\mathrm R}=0
\end{equation}
in the infrared and therefore
\begin{equation}
    \Gamma_{\mathrm D}=0
\end{equation}
for that realization.

Second, if the microscopic reference coupling does not project onto the
self-similar photon scale tangent with the overlap structure assumed in
Equation~\eqref{eq:AQ_eta}, the sixth-power suppression must be revised.

Third, if the observed fractional relaxation is strongly chromatic over a
frequency range in which finite-wave-number corrections are negligible, the
leading achromatic plateau closure is falsified.

Fourth, measurable cumulative spectral diffusion, image blurring,
polarization scrambling, photon destruction, or secondary-radiation
production at rates comparable to $\Gamma_{\mathrm D}$ would contradict the
proposed coherent scale-transfer mechanism.

Fifth, an independent determination of $r_{\mathrm D}$ significantly different
from the presently adopted value would modify $\Gamma_{\mathrm D}$ strongly
through
\begin{equation}
    \Gamma_{\mathrm D}
    \propto
    r_{\mathrm D}^{-7}.
\end{equation}

Sixth, if the full observational analysis requires identifying
$z_{\mathrm D}$ with the entire observed redshift but the resulting timing,
BAO, CMB, or other theory-independent observables cannot be reproduced, then
that \emph{cosmological identification} is falsified even if the microscopic
rate itself remains viable.

Seventh, if no complete spatial photon--receiver action can reproduce the
required real-frequency transport, spatial ray equation, and mechanical
recoil simultaneously, the corresponding microscopic implementation is
falsified.

Thus, failure of a complete cosmological interpretation does not
automatically imply that
\begin{equation}
    \Gamma_{\mathrm D}=0
\end{equation}
microscopically.

The microphysical mechanism and its macroscopic cosmological role are distinct
claims and must be tested separately.


\section{Summary}

The coherent photon branch of TFT requires a strict distinction between
ordinary amplitude damping and real-frequency scale transport.

The relevant coordinate is
\begin{equation}
    s=\ln q,
\end{equation}
with the desired local long-wave law
\begin{equation}
    \dot s
    =
    -\Gamma_{\mathrm D}.
\end{equation}

The established FCC transverse normalization is
\begin{equation}
    A_{\mathrm{ct}}^{(\mathrm{modal})}
    =
    2\pi\ell_{\mathrm S}^{2}.
\end{equation}

Under the adopted reference-mediated scale-transfer closure,
\begin{equation}
    V_{\mathrm R}
    =
    A_{\mathrm{ct}}^{(\mathrm{modal})}
    \ell_{\mathrm S}
    =
    2\pi\ell_{\mathrm S}^{3}.
\end{equation}

Relative to
\begin{equation}
    V_{\mathrm D}
    =
    \frac{4\pi}{3}
    r_{\mathrm D}^{3},
\end{equation}
the normalized overlap amplitude is
\begin{equation}
    \eta_{\mathrm R}
    =
    \sqrt{\frac32}
    \left(
        \frac{
            \ell_{\mathrm S}
        }{
            r_{\mathrm D}
        }
    \right)^{3/2}.
\end{equation}

Under the adopted two-reference-overlap-vertex return closure,
\begin{equation}
    P_{\mathrm R}
    =
    \eta_{\mathrm R}^{4}
    =
    \frac94
    \left(
        \frac{
            \ell_{\mathrm S}
        }{
            r_{\mathrm D}
        }
    \right)^6.
\end{equation}

Using
\begin{equation}
    t_{\mathrm{corr}}^{(0)}
    =
    \frac{
        4r_{\mathrm D}
    }{
        c_{\mathrm E}
    },
\end{equation}
the conditional microscopic scale-relaxation closure gives
\begin{equation}
    \Gamma_{\mathrm D}
    =
    \frac{9}{16}
    \mathcal F_{\mathrm R}
    \frac{
        c_{\mathrm E}
    }{
        r_{\mathrm D}
    }
    \left(
        \frac{
            \ell_{\mathrm S}
        }{
            r_{\mathrm D}
        }
    \right)^6.
\end{equation}

A finite nonzero $\Gamma_{\mathrm D}$ requires a carrier-projected
irreversible Active-Vacuum continuum. In the completed conservative coherent
sector,
\begin{equation}
    \mathcal F_{\mathrm R}=0,
    \qquad
    \Gamma_{\mathrm D}=0.
\end{equation}

For the unit-reference closure
\begin{equation}
    \mathcal F_{\mathrm R}=1,
\end{equation}
the presently adopted microscopic TFT scales give
\begin{equation}
    \Gamma_{\mathrm D}^{(0)}
    \simeq
    2.3591\times10^{-18}
    \ {\rm s^{-1}}.
\end{equation}

This numerical value is the result of the unit-reference
single-correlation closure. It is not a derivation of
$\mathcal F_{\mathrm R}=1$ from the complete nonlinear microscopic spectrum
and is not a physical lower bound on the relaxation rate.

Under preservation of the coherent modal action,
\begin{equation}
    I_{\gamma}
    =
    \frac{
        E_{\gamma}
    }{
        \omega_{\gamma}
    }
    =
    \hbar,
\end{equation}
the same conditional rate controls the real photon frequency:
\begin{equation}
    \frac{
        \dot E_{\gamma}
    }{
        E_{\gamma}
    }
    =
    \frac{
        \dot\omega_{\gamma}
    }{
        \omega_{\gamma}
    }
    =
    -\Gamma_{\mathrm D}.
\end{equation}

In the selected long-wave mean transport closure,
\begin{equation}
    \dot{\mathbf{k}}
    \simeq
    -\Gamma_{\mathrm D}\mathbf{k},
    \qquad
    \dot{\hat{\mathbf{k}}}=0,
    \qquad
    \dot N_{\mathrm{ph}}=0.
\end{equation}

The preservation of direction, polarization, photon number, and spectral
sharpness remains subject to the corresponding orthogonal-channel and
fluctuation constraints.

The energy removed from the coherent photon branch is transferred to
finite-energy noncoherent Active-Vacuum degrees of freedom. The Hard Floor is
not the dissipative receiving sector, and the present appendix does not
identify the transferred energy with a thermal state.

If the carrier-projected response possesses an approximately
$q$-independent plateau, the corresponding propagation-relaxation
frequency-shift factor is
\begin{equation}
    \ln(1+z_{\mathrm D})
    =
    \int_{\mathrm{path}}
    \frac{
        \Gamma_{\mathrm D}
    }{
        v_{\mathrm{ray}}
    }
    \,d\ell.
\end{equation}

This is not, by itself, the complete astronomical redshift.

In the stationary factorized regime,
\begin{equation}
    1+z_{\mathrm{obs}}
    =
    (1+z_{\mathrm{kin}})
    (1+z_{\mathrm{grav}})
    (1+z_{\mathrm D}).
\end{equation}

The corresponding Photon Energy-Relaxation Length,
\begin{equation}
    R_{\mathrm U}
    =
    \frac{
        c_{\mathrm E}
    }{
        \Gamma_{\mathrm D}
    },
\end{equation}
is the characteristic one-$e$-fold coherent photon scale-relaxation length
for a positive constant physical rate.

It is not a complete photon-energy-removal distance, a stopping boundary, or a
thermalization surface. It is also distinct from the Structural Reservoir
radius $R_{\mathrm{Res}}$, which characterizes the much larger macroscopic
structural extent of the TFT universe.

The microscopic scale-relaxation mechanism and the astronomical decomposition
of observed redshift are therefore deliberately separated. The present
appendix constructs a conditional microscopic closure for the former, while
the observational projection remains a distinct problem.
